%%%%%%%%%%%%%%%%%%%%%%%%%%%%%%%%%%%%%%%%
%
% $ Autor: Gruppe 07 $
% $ Projekt: Magic Faucet Fountain $
% $ Pfad: LaTeXTemplate/Nano33BLESense/Contets/General/Veränderungen/großeVeränderungen.tex $
% $ Kurzbeschreibung: Dieses Dokument zeigt die Liste der großen Veränderungen im Dokument. $
%
% !TeX spellcheck = de_DE/GB
% !TeX program = pdflatex
% !BIB program = biber/bibtex
% !TeX encoding = utf8
%
%%%%%%%%%%%%%%%%%%%%%%%%%%%%%%%%%%%%%%%%


\chapter{Große Änderungen in \menu{LaTeXTemplate}:}
\textbf{Große Änderungen sind:} Dokument umgeschrieben/übersetzt/neu erstellt


\section{\menu{CAD und 3D-Druck}}
	\begin{itemize}
	\item \menu{CAD Dateien} erstellt
		\begin{itemize}
		\item CAD-Dateien der konstruierten Teile eingefügt
	    \end{itemize}
	\item \menu{GCode} erstellt
		\begin{itemize}
		\item G-Code der 3D-Druck-Teile eingefügt
	    \end{itemize}
    \end{itemize}
    
    
\section{\menu{Demontageanleitung}}
\begin{itemize}
	\item \menu{General}
	    \begin{itemize}
		\item Abbildung \python{blumentopf} eingefügt
    	\end{itemize}
	\item \python{Kapitel 1: Einleitung}
	    \begin{itemize}
		\item Einleitung zu der Demontage geschrieben
      	\end{itemize}
	\item \python{Kapitel 2: WerkzeugeHilfsmittel}
		\begin{itemize}
		\item Erfolderlichen Werkzeuge und Hilfmittel aufgezählt und erläutert
     	\end{itemize}
	\item \python{Kapitel 3: Demontageablauf}
		\begin{itemize}
		\item Demontageprozess beschrieben
		\item Tikzbilder zur Verdeutlichung des Demontageprozesses erstellt
	    \end{itemize}
	\item \python{Kapitel 4: Lagerung}
		\begin{itemize}
		\item Empfehlungen zur Lagerung der Baugruppe geschrieben
    	\end{itemize}
\end{itemize}


\section{\menu{Handbuch Magic Faucet Fountain}}
\begin{itemize}
	\item \menu{General}
	    \begin{itemize}
		\item Abbildungen \python{netzteilHandbuch}, \python{steine} und \python{topfHandbuch} eingefügt
    	\end{itemize}
	\item \python{Kapitel 1: Hinweise}
    	\begin{itemize}
		\item Bedienungshinweise und Sicherheitshinweise geschrieben
		\item Danksagung für den Kauf des Produkts erstellt
    	\end{itemize}
	\item \python{Kapitel 2: Hauptfunktion}
    	\begin{itemize}
		\item Text zu der Hauptfunktion geschrieben
		\item Erste Schritte zur Anwendung beschrieben
		\item Hinweise zur Aufbau der Baugruppe beschrieben
		\item Einrichten der App beschrieben
		\item Tikzbilder der Baugruppe und des Netzteilkastens erstellt
    	\end{itemize}
	\item \python{Kapitel 3: Skizze}
    	\begin{itemize}
		\item Tikzbild der Übersichtsskizze mit Beschriftung erstellt
    	\end{itemize}
	\item \python{Kapitel 4: Liste der Teile}
    	\begin{itemize}
		\item Text zu den verwendetet Bauteilen geschrieben
		\item Tabelle mit verwendeten Teilen erstellt
		\item Bilder der verwendeten Teile in die Tabelle eingefügt
    	\end{itemize}
	\item \python{Kapitel 5: Aufbau}
	    \begin{itemize}
		\item Generellen Aufbau der Baugruppe geschrieben
		\item Wasserführung der Baugruppe beschrieben
		\item Pegelmessung der Baugruppe beschrieben
		\item Elektronik und Steuerung der Baugruppe erklärt
    	\end{itemize}
	\item \python{Kapitel 6: Funktionen}
		\begin{itemize}
		\item Funktionen der Baugruppe beschrieben
	    \end{itemize}
	\item \python{Kapitel 7: Wartung}
		\begin{itemize}
		\item Allgemeine Hinweise und Wartungshinweise in Bezug auf die Baugruppe geschrieben
	    \end{itemize}
	\item \python{Kapitel 8: Fehlerbehandlung}
		\begin{itemize}
		\item Tabelle mit möglichen Problemen und den dazugehörigen Behebungsmöglichkeiten
    	\end{itemize}
\end{itemize}


\section{\menu{MLBib}}
\textbf{Änderungen in \menu{MLBib}}
\begin{itemize}
	\item \menu{PDF Literatur} erstellt
	\begin{itemize}
		\item PDF-Dateien der verwendeten Literatur eingefügt
	\end{itemize}
	\item \menu{Präsentation der Literatur} erstellt
		\begin{itemize}
		\item Eine Präsentation der verwendeten Literatur mithilfe des LaTeX-Templates erstellt
		\item Bilder der verwendetet Quellen eingefügt
	\end{itemize}
\end{itemize}


\section{\menu{Montageanleitung}}
\begin{itemize}
	\item \menu{General}
    	\begin{itemize}
		\item Abbildung \python{topfMontage} eingefügt
    	\end{itemize}
	\item \python{Kapitel 1: Einleitung}
    	\begin{itemize}
		\item Einleitung zu der Montage geschrieben
		\item Hinweis zu der richtigen Anwendung der Baugruppe erstellt
    	\end{itemize}
	\item \python{Kapitel 2: Material- und Werkzeugliste}
    	\begin{itemize}
		\item Erfolderlichen Werkzeuge und Materialien aufgezählt und erläutert
		\item Tikzbilder der Baugruppe erstellt
		\item Tabelle der Materialien erstellt
    	\end{itemize}
	\item \python{Kapitel 3: Montageablauf}
    	\begin{itemize}
		\item Montageprozess beschrieben
		\item Tikzbilder zur Verdeutlichung des Montageprozesses erstellt
     	\end{itemize}
	\item \python{Kapitel 4: Inbetriebnahme und Funktionstest}
	    \begin{itemize}
		\item Die Inbetriebnahme der Baugruppe erklärt
	    \item Funktionstest beschrieben
	    \end{itemize}
	\item \python{Kapitel 5: Abschluss}
		\begin{itemize}
		\item Hinweis zum Abschluss der Arbeit geschrieben
    	\end{itemize}
\end{itemize}

	
	
\section{\menu{Nano33BLESense}}
\subsection{\menu{Nano33BLESense>Code}}
 \textbf{Änderungen in \menu{Nano33BLESense > Code > Arduino}}
\begin{itemize}
	\item \menu{Applikation>ablaufApplikation} erstellt
	\item \menu{Bluetooth Modul>bluetoothAnwendung} erstellt
	\item \menu{Bluetooth Modul>bluetoothBeispiel} erstellt
	\item \menu{Externe LED>ansteuerungLED} erstellt
	\item \menu{Externe LED>blinkenLED} erstellt
	\item \menu{MOSFET} erstellt
	\item \menu{NinaB306>anwendungBLE} erstellt
	\item \menu{NinaB306>testBLE} erstellt
	\item \menu{nRF Connect>einrichtungNRF} erstellt
	\item \menu{Pumpe>ansteuerungPumpe} erstellt
	\item \menu{Pumpe>fehlererkennungPumpe} erstellt
	\item \menu{Relais>anwendungRelais} erstellt
	\item \menu{Relais>einfachRelaiss} erstellt
\end{itemize}
 \textbf{Änderungen in \menu{Nano33BLESense > Contents > General}}
 \begin{itemize}
 	\item \menu{HardwareBeschreibung} erstellt
 	\item \menu{Schaltplan} erstellt
 	\item \menu{Veränderungen} erstellt
 \end{itemize}
Einleitung Hauptdokument neu erstellt und geschrieben: \ref{chap:einleitungDoku}

	
\section{\menu{Poster}}
    \begin{itemize}
    \item \menu{images}
	    \begin{itemize}
		\item Abbildung \python{topfPoster} eingefügt
		\item Abbildung \python{Rendering8} entfernt
		\item Abbildung \python{Rendering9} entfernt
		\item Abbildung \python{Rendering10} entfernt
    	\end{itemize}
    \item \python{PosterMagicFaucetFountain}
        \begin{itemize}
     	\item Poster mithilfe des Templates erstellt
     	\item Bild der Baugruppe eingefügt
        \end{itemize}
	\end{itemize}
		
\section{\menu{Präsentation}}


\section{\menu{Schnelleinstieg}}

	
\section{\menu{Xmind}}
	
	
\section{\menu{author}}
	
	
\section{\menu{README}}

